
\ctikzset{%
    % voltage/european label distance=0.5
    % Colors
    amplifiers/fill=iceberg,
    capacitors/fill=palesilver,
    resistors/fill=meatbrown,
    %sources/fill=turquoise,
    %bipoles/esource/voltage/distance from node/.initial=0.1,
    %bipoles/capacitor/voltage/distance from node/.initial=0.0,
    voltage/distance from node/.initial=0.2,
    voltage shift=1.0,
    bipoles/cuteswitch/thickness=0.4,
    bipoles/cuteswitch/shape=circ,
    bipoles/cuteswitch/height=0.5,
    %bipoles/cuteswitch/width=0.4,
    %misc/fill=xanadu,
    % From CircuiTikZ
}
\tikzset{
    longL/.style={%
        cute inductor,inductors/scale=1.2,
        inductors/width=1.6,inductors/coils=11
    }
}

% From the manual
\NewDocumentCommand{\fixedvlen}{O{0.5cm} m m O{}}{% [semilength]{node}{label}[extra options]
    % get the center of the standard arrow
    \coordinate (#2-Vcenter) at ($(#2-Vfrom)!0.5!(#2-Vto)$);
    % draw an arrow of a fixed size around that center and on the same line
    \draw[-Triangle, #4] ($(#2-Vcenter)!#1!(#2-Vfrom)$) -- ($(#2-Vcenter)!#1!(#2-Vto)$);
    % position the label as in the normal voltages
    \node[anchor=\ctikzgetanchor{#2}{Vlab}, #4] at (#2-Vlab) {#3};
}
